\import{cardinal/exponentiation.tex}
\import{nat/multiplication.tex}
\import{set/finite.tex}
\import{set/partition.tex}

\section{Finite and countable cardinality}

The cardinal representative agrees with the von Neumann numeral on finite
sets and with $\naturals$ on countably infinite sets.

\begin{proposition}\label{successor_remove_element_equinum}
    Suppose $y\in\suc{n}$.
    Then $\suc{n}\setminus\{y\}\approx n$.
\end{proposition}
\begin{proof}
    \begin{byCase}
        \caseOf{$y=n$.}
            $\suc{n}\setminus\{n\}=n$
                by \auto{setext,setminus,singleton_iff,suc_iff,in_irrefl}.
            Follows by \obv{equinum_refl}.
        \caseOf{$y\neq n$.}
            Hence $y\in n$ by \auto{suc_iff}.
            Let $A=n\setminus\{y\}$.
            $A\inter\{n\}=\emptyset$
                by
                \auto{setext,inter,setminus,singleton_iff,emptyset,in_irrefl}.
            $A\inter\{y\}=\emptyset$
                by \auto{setext,inter,setminus,singleton_iff,emptyset}.
            $A\approx A$ by \obv{equinum_refl}.
            $\{n\}\approx\{y\}$ by \obv{singletons_equinum}.
            Hence $A\union\{n\}\approx A\union\{y\}$
                by \obv{disjoint_union_equinum_congruent}.
            $A\union\{n\}=\suc{n}\setminus\{y\}$
                by
                \auto{setext,union_iff,setminus,singleton_iff,suc_iff,in_irrefl}.
            $A\union\{y\}=n$
                by
                \auto{setext,union_iff,setminus,singleton_iff}.
    \end{byCase}
\end{proof}

\begin{theorem}\label{natural_equinumerous_unique}
    Let $m,n$ be natural numbers.
    Suppose $m\approx n$.
    Then $m=n$.
\end{theorem}
\begin{proof}
    Let $C=\{k\in\naturals\mid
        \text{for every $j\in\naturals$ such that $j\approx k$
            we have $j=k$}\}$.

    Show that $\emptyset\in C$.
    \begin{subproof}
        $\emptyset\in\naturals$ by \obv{omega_empty}.
        Show that for every $j\in\naturals$ such that
            $j\approx\emptyset$ we have $j=\emptyset$.
        \begin{subproof}
            Fix $j$.
            Assume $j\in\naturals$ and $j\approx\emptyset$.
            Take $f$ such that $f\in\bijections{j}{\emptyset}$
                by \obv{equinum}.
            Suppose not.
            Hence $j$ is inhabited by \obv{nonempty_inhabited}.
            Take $x$ such that $x\in j$.
            $f(x)\in\emptyset$
                by \auto{bijections_to_funs,funs_type_apply}.
            Contradiction by \obv{emptyset}.
        \end{subproof}
    \end{subproof}

    Show that for every $k\in C$ we have $\suc{k}\in C$.
    \begin{subproof}
        Fix $k\in C$.
        $\suc{k}$ is a natural number by \obv{omega_successor}.
        Show that for every $j\in\naturals$ such that
            $j\approx\suc{k}$ we have $j=\suc{k}$.
        \begin{subproof}
            Fix $j$.
            Assume $j\in\naturals$ and $j\approx\suc{k}$.
            Take $f$ such that $f\in\bijections{j}{\suc{k}}$
                by \obv{equinum}.
            Show that $j\neq\emptyset$.
            \begin{subproof}
                Suppose not.
                $\ran{f}=\suc{k}$ by \auto{bijections,ran_of_surj}.
                Hence $k\in\ran{f}$ by \obv{suc_intro_self}.
                Take $x$ such that $(x,k)\in f$ by \obv{ran_iff}.
                $x\in\dom{f}$ by \obv{dom_intro}.
                $\dom{f}=j$ by \obv{bijections_dom}.
                Hence $x\in\emptyset$.
                Contradiction by \obv{emptyset}.
            \end{subproof}
            Take a natural number $l$ such that $j=\suc{l}$
                by \obv{natural_zero_or_successor}.
            Hence $f\in\bijections{\suc{l}}{\suc{k}}$.
            Let $z=f(l)$.
            $l\in\suc{l}$ by \obv{suc_intro_self}.
            Hence $z\in\suc{k}$
                by \auto{bijections_to_funs,funs_type_apply}.
            $\restrl{f}{l}\in
                \bijections{l}{\suc{k}\setminus\{z\}}$
                by \auto{bijection_remove_successor}.
            Hence $l\approx\suc{k}\setminus\{z\}$ by \obv{equinum}.
            $\suc{k}\setminus\{z\}\approx k$
                by \obv{successor_remove_element_equinum}.
            Hence $l\approx k$ by \obv{equinum_tran}.
            Thus $l=k$.
        \end{subproof}
    \end{subproof}

    Hence $C$ is an inductive set by \obv{inductive_set}.
    Hence every natural number is an element of $C$
        by \obv{omega_minimal}.
    Thus $n\in C$.
\end{proof}

\subsection{Finite-cardinality witnesses}

The following counting lemmas use explicit finite bijections only.  In
particular, they do not pass through canonical cardinals or choose a family of
representatives.

\begin{proposition}\label{finite_cardinality_bijection}
    Suppose $X$ has finite cardinality $n$ and
        $f\in\bijections{X}{Y}$.
    Then $Y$ has finite cardinality $n$.
\end{proposition}
\begin{proof}
    Take $g$ such that $g\in\bijections{n}{X}$ by
        \obv{finite_cardinality}.
    $f\circ g\in\bijections{n}{Y}$ by \obv{bijection_circ}.
    Follows by \obv{finite_cardinality}.
\end{proof}

\begin{theorem}\label{finite_cardinality_unique}
    Suppose $X$ has finite cardinality $m$ and
        $X$ has finite cardinality $n$.
    Then $m=n$.
\end{theorem}
\begin{proof}
    $m,n\in\naturals$ by \obv{finite_cardinality}.
    Take $f$ such that $f\in\bijections{m}{X}$ by
        \obv{finite_cardinality}.
    Take $g$ such that $g\in\bijections{n}{X}$ by
        \obv{finite_cardinality}.
    $\converse{g}\in\bijections{X}{n}$ by
        \obv{bijection_converse_is_bijection}.
    Hence $\converse{g}\circ f\in\bijections{m}{n}$ by
        \obv{bijection_circ}.
    Hence $m\approx n$ by \obv{equinum}.
    Follows by \obv{natural_equinumerous_unique}.
\end{proof}

\begin{proposition}\label{finite_cardinality_zero_iff_emptyset}
    $X$ has finite cardinality $\emptyset$ iff $X=\emptyset$.
\end{proposition}
\begin{proof}
    Show that if $X$ has finite cardinality $\emptyset$, then
        $X=\emptyset$.
    \begin{subproof}
        Assume $X$ has finite cardinality $\emptyset$.
        Take $f$ such that $f\in\bijections{\emptyset}{X}$ by
            \obv{finite_cardinality}.
        $\ran{f}=X$ by \auto{bijections,ran_of_surj}.
        $\dom{f}=\emptyset$ by \obv{bijections_dom}.
        Hence $f=\emptyset$ by
            \auto{setext,emptyset,bijection_is_function,
                function_member_elim}.
        Follows by \obv{ran_emptyset}.
    \end{subproof}
    Show that if $X=\emptyset$, then
        $X$ has finite cardinality $\emptyset$.
    \begin{subproof}
        Assume $X=\emptyset$.
        $\identity{\emptyset}\in
            \bijections{\emptyset}{\emptyset}$ by
            \obv{id_elem_bijections}.
        Follows by \obv{finite_cardinality,omega_empty}.
    \end{subproof}
    Follows by \auto.
\end{proof}

\begin{proposition}\label{finite_cardinality_cons}
    Let $n$ be a natural number.
    Suppose $X$ has finite cardinality $n$ and $x\notin X$.
    Then $\cons{x}{X}$ has finite cardinality $\suc{n}$.
\end{proposition}
\begin{proof}
    Take $f$ such that $f\in\bijections{n}{X}$ by
        \obv{finite_cardinality}.
    $n\notin n$ by \obv{in_irrefl}.
    $\cons{(n,x)}{f}\in
        \bijections{\cons{n}{n}}{\cons{x}{X}}$ by
        \obv{bijection_cons}.
    $\cons{n}{n}=\suc{n}$ by \obv{suc,cons_iff,setext}.
    $\suc{n}$ is a natural number by \obv{omega_successor}.
    Follows by \auto{finite_cardinality}.
\end{proof}

\begin{proposition}\label{finite_cardinality_remove_element}
    Let $n$ be a natural number.
    Suppose $X$ has finite cardinality $\suc{n}$ and $x\in X$.
    Then $X\setminus\{x\}$ has finite cardinality $n$.
\end{proposition}
\begin{proof}
    Take $f$ such that $f\in\bijections{\suc{n}}{X}$ by
        \obv{finite_cardinality}.
    Take $j\in\suc{n}$ such that $f(j)=x$ by
        \auto{bijections,surj}.
    Let $D=\suc{n}\setminus\{j\}$ and $g=\restrl{f}{D}$.
    $f\in\funs{\suc{n}}{X}$ by \obv{bijections_to_funs}.
    Hence $f$ is a function and $\dom{f}=\suc{n}$ by
        \obv{funs_is_function,funs}.
    $f$ is injective by \obv{bijections}.
    Hence $f\in\inj{\suc{n}}{X}$ by
        \auto{inj,injective_function}.
    $D\subseteq\suc{n}$ by \auto{setminus_subseteq}.
    Hence $g\in\inj{D}{X}$ by \auto{inj_restrl}.
    Show that $\ran{g}=X\setminus\{x\}$.
    \begin{subproof}
        Show that every element of $\ran{g}$ is an element of
            $X\setminus\{x\}$.
        \begin{subproof}
            Fix $y\in\ran{g}$.
            Take $i$ such that $(i,y)\in g$ by \obv{ran_iff}.
            Hence $i\in D$ and $(i,y)\in f$ by \auto{restrl_iff}.
            Hence $y\in\ran{f}$ by \obv{ran_intro}.
            $\ran{f}=X$ by \auto{bijections,ran_of_surj}.
            Hence $y\in X$.
            $i\neq j$ by \auto{setminus,singleton_iff}.
            Show that $y\neq x$.
            \begin{subproof}
                Suppose not.
                $j\in\dom{f}$.
                $(j,f(j))\in f$ by \obv{function_apply_elim}.
                Hence $(j,x)\in f$.
                Hence $i=j$ by \auto{injective}.
                Contradiction.
            \end{subproof}
            Follows by \auto{setminus,singleton_iff}.
        \end{subproof}
        Show that every element of $X\setminus\{x\}$ is an element of
            $\ran{g}$.
        \begin{subproof}
            Fix $y\in X\setminus\{x\}$.
            Hence $y\in X$ and $y\neq x$ by
                \auto{setminus,singleton_iff}.
            Take $i\in\suc{n}$ such that $f(i)=y$ by
                \auto{bijections,surj}.
            $i\neq j$.
            Hence $i\in D$ by \auto{setminus,singleton_iff}.
            $(i,y)\in f$ by \auto{function_apply_elim,bijection_is_function,
                bijections_dom}.
            Hence $(i,y)\in g$ by \obv{restrl_iff}.
            Follows by \obv{ran_intro}.
        \end{subproof}
        Follows by set extensionality.
    \end{subproof}
    $g\in\surj{D}{\ran{g}}$ by \auto{inj,fun_to_surj}.
    Hence $g\in\surj{D}{X\setminus\{x\}}$.
    $g$ is injective by \obv{restrl_injective}.
    Hence $g\in\bijections{D}{X\setminus\{x\}}$ by
        \obv{bijections}.
    $D\approx n$ by \auto{successor_remove_element_equinum}.
    Hence $n\approx D$ by \obv{equinum_sym}.
    Take $h$ such that $h\in\bijections{n}{D}$ by \obv{equinum}.
    $g\circ h\in\bijections{n}{X\setminus\{x\}}$ by
        \obv{bijection_circ}.
    Follows by \obv{finite_cardinality}.
\end{proof}

\begin{theorem}\label{finite_cardinality_disjoint_union}
    Let $m,n$ be natural numbers.
    Suppose $A$ has finite cardinality $m$ and
        $B$ has finite cardinality $n$ and $A\inter B=\emptyset$.
    Then $A\union B$ has finite cardinality $m+n$.
\end{theorem}
\begin{proof}
    Let $C=\{k\in\naturals\mid
        \text{for all $Y$ such that $Y$ has finite cardinality $k$ and
            $A\inter Y=\emptyset$ we have
            $A\union Y$ has finite cardinality $m+k$}\}$.
    Show that $\emptyset\in C$.
    \begin{subproof}
        $\emptyset\in\naturals$ by \obv{omega_empty}.
        Show that for all $Y$ such that
            $Y$ has finite cardinality $\emptyset$ and
            $A\inter Y=\emptyset$ we have
            $A\union Y$ has finite cardinality $m+\emptyset$.
        \begin{subproof}
            Fix $Y$.
            Assume $Y$ has finite cardinality $\emptyset$ and
                $A\inter Y=\emptyset$.
            $Y=\emptyset$ by \obv{finite_cardinality_zero_iff_emptyset}.
            $A\union Y=A$ by \auto{union_emptyset}.
            $m+\emptyset=m$ by \obv{num_addition_zero_right}.
        \end{subproof}
    \end{subproof}
    Show that for every $k\in C$ we have $\suc{k}\in C$.
    \begin{subproof}
        Fix $k\in C$.
        Show that for all $Y$ such that
            $Y$ has finite cardinality $\suc{k}$ and
            $A\inter Y=\emptyset$ we have
            $A\union Y$ has finite cardinality $m+\suc{k}$.
        \begin{subproof}
            Fix $Y$.
            Assume $Y$ has finite cardinality $\suc{k}$ and
                $A\inter Y=\emptyset$.
            Take $f$ such that $f\in\bijections{\suc{k}}{Y}$ by
                \obv{finite_cardinality}.
            Let $y=f(k)$ and $Y'=Y\setminus\{y\}$.
            $k\in\suc{k}$ by \obv{suc_intro_self}.
            Hence $y\in Y$ by
                \auto{bijections_to_funs,funs_type_apply}.
            $Y'$ has finite cardinality $k$ by
                \auto{finite_cardinality_remove_element}.
            $A\inter Y'=\emptyset$ by
                \auto{setext,inter,setminus,emptyset}.
            Hence $A\union Y'$ has finite cardinality $m+k$.
            $y\notin A\union Y'$ by
                \auto{union_iff,setminus,singleton_iff,inter,emptyset}.
            Hence $\cons{y}{A\union Y'}$ has finite cardinality
                $\suc{m+k}$ by
                \obv{finite_cardinality_cons,num_addition_on_naturals}.
            $\cons{y}{A\union Y'}=A\union Y$ by
                \auto{setext,cons_iff,union_iff,setminus,singleton_iff}.
            $m+\suc{k}=\suc{m+k}$ by
                \obv{num_addition_successor_right}.
        \end{subproof}
        Follows by \obv{omega_successor}.
    \end{subproof}
    Hence $C$ is an inductive set by \obv{inductive_set}.
    Hence every natural number is an element of $C$ by
        \obv{omega_minimal}.
    Thus $n\in C$.
\end{proof}

\begin{theorem}\label{finite_subset_same_cardinality}
    Let $n$ be a natural number.
    Suppose $A\subseteq B$ and $A$ and $B$ have finite cardinality $n$.
    Then $A=B$.
\end{theorem}
\begin{proof}
    Let $C=\{k\in\naturals\mid
        \text{for all $X,Y$ such that $X\subseteq Y$ and
            $X$ and $Y$ have finite cardinality $k$ we have $X=Y$}\}$.
    $\emptyset\in C$ by
        \auto{omega_empty,finite_cardinality_zero_iff_emptyset}.
    Show that for every $k\in C$ we have $\suc{k}\in C$.
    \begin{subproof}
        Fix $k\in C$.
        Show that for all $X,Y$ such that $X\subseteq Y$ and
            $X$ and $Y$ have finite cardinality $\suc{k}$ we have
            $X=Y$.
        \begin{subproof}
            Fix $X,Y$.
            Assume $X\subseteq Y$ and $X$ and $Y$ have finite cardinality
                $\suc{k}$.
            Take $f$ such that $f\in\bijections{\suc{k}}{X}$ by
                \obv{finite_cardinality}.
            Let $x=f(k)$.
            $k\in\suc{k}$ by \obv{suc_intro_self}.
            Hence $x\in X$ by
                \auto{bijections_to_funs,funs_type_apply}.
            Hence $x\in Y$ by \obv{subseteq}.
            $X\setminus\{x\}$ and $Y\setminus\{x\}$ have finite
                cardinality $k$ by
                \obv{finite_cardinality_remove_element}.
            $X\setminus\{x\}\subseteq Y\setminus\{x\}$ by
                \obv{subseteq,setminus}.
            Hence $X\setminus\{x\}=Y\setminus\{x\}$.
            $\cons{x}{X\setminus\{x\}}=X$ by \obv{cons_remove}.
            $\cons{x}{Y\setminus\{x\}}=Y$ by \obv{cons_remove}.
            \begin{align*}
                X
                    &= \cons{x}{X\setminus\{x\}}
                    \\
                    &= \cons{x}{Y\setminus\{x\}}
                    \\
                    &= Y
            \end{align*}
        \end{subproof}
        Follows by \obv{omega_successor}.
    \end{subproof}
    Hence $C$ is an inductive set by \obv{inductive_set}.
    Hence every natural number is an element of $C$ by
        \obv{omega_minimal}.
    Thus $n\in C$.
\end{proof}

\begin{theorem}\label{uniform_finite_partition_cardinality}
    Let $n,k$ be natural numbers.
    Suppose $P$ is a partition of $X$.
    Suppose $P$ has finite cardinality $k$ and
        every element of $P$ has finite cardinality $n$.
    Then $X$ has finite cardinality $\naturalMul{n}{k}$.
\end{theorem}
\begin{proof}
    Let $C=\{l\in\naturals\mid
        \text{for all $Y,Q$ such that $Q$ is a partition of $Y$ and
            $Q$ has finite cardinality $l$ and every element of $Q$
            has finite cardinality $n$ we have
            $Y$ has finite cardinality $\naturalMul{n}{l}$}\}$.
    Show that $\emptyset\in C$.
    \begin{subproof}
        $\emptyset\in\naturals$ by \obv{omega_empty}.
        Show that for all $Y,Q$ such that $Q$ is a partition of $Y$ and
            $Q$ has finite cardinality $\emptyset$ and every element
            of $Q$ has finite cardinality $n$ we have
            $Y$ has finite cardinality $\naturalMul{n}{\emptyset}$.
        \begin{subproof}
            Fix $Y,Q$.
            Assume $Q$ is a partition of $Y$ and
                $Q$ has finite cardinality $\emptyset$ and every
                element of $Q$ has finite cardinality $n$.
            $Q=\emptyset$ by
                \obv{finite_cardinality_zero_iff_emptyset}.
            $\unions{Q}=Y$.
            $\unions{Q}=\emptyset$ by \obv{unions_emptyset}.
            Hence $Y=\emptyset$.
            $\naturalMul{n}{\emptyset}=\emptyset$ by
                \obv{natural_multiplication_zero_right}.
            Follows by \obv{finite_cardinality_zero_iff_emptyset}.
        \end{subproof}
    \end{subproof}
    Show that for every $l\in C$ we have $\suc{l}\in C$.
    \begin{subproof}
        Fix $l\in C$.
        Show that for all $Y,Q$ such that $Q$ is a partition of $Y$ and
            $Q$ has finite cardinality $\suc{l}$ and every element
            of $Q$ has finite cardinality $n$ we have
            $Y$ has finite cardinality $\naturalMul{n}{\suc{l}}$.
        \begin{subproof}
            Fix $Y,Q$.
            Assume $Q$ is a partition of $Y$ and
                $Q$ has finite cardinality $\suc{l}$ and every
                element of $Q$ has finite cardinality $n$.
            Take $f$ such that $f\in\bijections{\suc{l}}{Q}$ by
                \obv{finite_cardinality}.
            Let $B=f(l)$, $Q'=Q\setminus\{B\}$, and $Y'=Y\setminus B$.
            $l\in\suc{l}$ by \obv{suc_intro_self}.
            Hence $B\in Q$ by
                \auto{bijections_to_funs,funs_type_apply}.
            $Q'$ is a partition of $Y'$ by
                \auto{partition_remove_block}.
            $Q'$ has finite cardinality $l$ by
                \auto{finite_cardinality_remove_element}.
            Every element of $Q'$ has finite cardinality $n$ by
                \auto{setminus}.
            Hence $Y'$ has finite cardinality $\naturalMul{n}{l}$.
            $B$ has finite cardinality $n$.
            $Y'\inter B=\emptyset$ by
                \auto{setext,inter,setminus,emptyset}.
            Hence $Y'\union B$ has finite cardinality
                $\naturalMul{n}{l}+n$ by
                \obv{finite_cardinality_disjoint_union,
                    natural_multiplication_closed}.
            $B\subseteq Y$ by \obv{partition_block_subseteq}.
            Hence $Y'\union B=Y$ by
                \auto{setext,union_iff,setminus,subseteq}.
            $\naturalMul{n}{\suc{l}}=\naturalMul{n}{l}+n$ by
                \obv{natural_multiplication_successor_right}.
        \end{subproof}
        Follows by \obv{omega_successor}.
    \end{subproof}
    Hence $C$ is an inductive set by \obv{inductive_set}.
    Hence every natural number is an element of $C$ by
        \obv{omega_minimal}.
    Thus $k\in C$.
\end{proof}

\begin{theorem}\label{natural_number_is_cardinal}
    Let $n$ be a natural number.
    Then $n$ is a cardinal.
\end{theorem}
\begin{proof}
    $n$ is an ordinal by \obv{natural_number_is_ordinal}.
    Show that for every $m\in n$ we have $\lnot(m\approx n)$.
    \begin{subproof}
        Fix $m\in n$.
        $n\subseteq\naturals$
            by
            \obv{naturals_is_transitive_set,transitiveset_iff_subseteq}.
        Suppose $m\approx n$.
        Hence $m=n$ by \obv{natural_equinumerous_unique}.
        Contradiction by \obv{in_irrefl}.
    \end{subproof}
    Hence $n$ is an initial ordinal by \obv{initial_ordinal}.
    Follows by \obv{cardinal}.
\end{proof}

\begin{corollary}\label{card_of_natural_number}
    Let $n$ be a natural number.
    Then $\card{n}=n$.
\end{corollary}
\begin{proof}
    Follows by \obv{natural_number_is_cardinal,card_of_cardinal}.
\end{proof}

\begin{theorem}\label{finite_equinumerous}
    Suppose $X$ is finite and $X\approx Y$.
    Then $Y$ is finite.
\end{theorem}
\begin{proof}
    Take a natural number $n$ such that there exists
        $f\in\bijections{n}{X}$ by \auto{finite}.
    Take $f$ such that $f\in\bijections{n}{X}$ by \auto.
    Hence $n\approx X$ by \obv{equinum}.
    Hence $n\approx Y$ by \obv{equinum_tran}.
    Take $g$ such that $g\in\bijections{n}{Y}$ by \obv{equinum}.
    Follows by \obv{finite}.
\end{proof}

\begin{theorem}\label{finite_cardinality_iff_card}
    Let $n$ be a natural number.
    Then $X$ has finite cardinality $n$ iff $\card{X}=n$.
\end{theorem}
\begin{proof}
    Show that if $X$ has finite cardinality $n$, then $\card{X}=n$.
    \begin{subproof}
        Assume $X$ has finite cardinality $n$.
        Take $f$ such that $f\in\bijections{n}{X}$
            by \obv{finite_cardinality}.
        Hence $n\approx X$ by \obv{equinum}.
        Hence $\card{n}=\card{X}$
            by \obv{card_eq_iff_equinumerous}.
        $\card{n}=n$ by \obv{card_of_natural_number}.
    \end{subproof}
    Show that if $\card{X}=n$, then
        $X$ has finite cardinality $n$.
    \begin{subproof}
        Assume $\card{X}=n$.
        $X\approx\card{X}$ by \obv{equinumerous_card}.
        Hence $X\approx n$.
        Hence $n\approx X$ by \obv{equinum_sym}.
        Take $f$ such that $f\in\bijections{n}{X}$ by \obv{equinum}.
        Follows by \obv{finite_cardinality}.
    \end{subproof}
    Follows by \auto.
\end{proof}

\begin{proposition}\label{finite_cardinality_one_elements_equal}
    Suppose $X$ has finite cardinality $\cardOne$.
    Let $x,y\in X$.
    Then $x=y$.
\end{proposition}
\begin{proof}
    Take $f$ such that $f\in\bijections{\cardOne}{X}$ by
        \obv{finite_cardinality}.
    Take $i,j\in\cardOne$ such that $f(i)=x$ and $f(j)=y$ by
        \auto{bijections,surj}.
    $i=\emptyset$ and $j=\emptyset$ by
        \auto{cardinal_one,suc_iff,emptyset}.
\end{proof}

\begin{proposition}\label{finite_cardinality_one_iff_singleton}
    $X$ has finite cardinality $\cardOne$ iff
        there exists $x$ such that $X=\{x\}$.
\end{proposition}
\begin{proof}
    $\cardOne$ is a natural number by
        \auto{cardinal_one,omega_empty,omega_successor}.
    Show that if $X$ has finite cardinality $\cardOne$, then
        there exists $x$ such that $X=\{x\}$.
    \begin{subproof}
        Assume $X$ has finite cardinality $\cardOne$.
        Take $f$ such that $f\in\bijections{\cardOne}{X}$ by
            \obv{finite_cardinality}.
        $\emptyset\in\cardOne$ by
            \obv{cardinal_one,suc_intro_self}.
        Let $x=f(\emptyset)$.
        $x\in X$ by \auto{bijections_to_funs,funs_type_apply}.
        Show that every element of $X$ is equal to $x$.
        \begin{subproof}
            Follows by \obv{finite_cardinality_one_elements_equal}.
        \end{subproof}
        Hence $X=\{x\}$ by \auto{setext,singleton_iff}.
    \end{subproof}
    Show that if there exists $x$ such that $X=\{x\}$, then
        $X$ has finite cardinality $\cardOne$.
    \begin{subproof}
        Assume there exists $x$ such that $X=\{x\}$.
        Take $x$ such that $X=\{x\}$.
        $\cardOne=\{\emptyset\}$ by
            \obv{cardinal_one,setext,suc_iff,singleton_iff,emptyset}.
        $\{\emptyset\}\approx\{x\}$ by \obv{singletons_equinum}.
        Hence $\cardOne\approx X$.
        Take $f$ such that $f\in\bijections{\cardOne}{X}$ by
            \obv{equinum}.
        Follows by \obv{finite_cardinality}.
    \end{subproof}
    Follows by \auto.
\end{proof}

\begin{theorem}\label{two_element_cardinality_characterization}
    $A\approx\cardTwo$ iff there exist $x,y$ such that
        $x\neq y$ and $A=\{x,y\}$.
\end{theorem}
\begin{proof}
    $\cardOne\in\naturals$ by
        \auto{cardinal_one,omega_empty,omega_successor}.
    Hence $\cardTwo\in\naturals$ by
        \auto{cardinal_two_successor,omega_successor}.
    Show that if $A\approx\cardTwo$, then there exist $x,y$ such that
        $x\neq y$ and $A=\{x,y\}$.
    \begin{subproof}
        Assume $A\approx\cardTwo$.
        Hence $\cardTwo\approx A$ by \obv{equinum_sym}.
        Take $f$ such that $f\in\bijections{\cardTwo}{A}$ by
            \obv{equinum}.
        Hence $A$ has finite cardinality $\cardTwo$ by
            \obv{finite_cardinality}.
        $\cardOne\in\cardTwo$
            by \obv{cardinal_two_successor,suc_intro_self}.
        Let $y=f(\cardOne)$.
        $y\in A$ by \auto{bijections_to_funs,funs_type_apply}.
        $A\setminus\{y\}$ has finite cardinality $\cardOne$ by
            \obv{cardinal_two_successor,finite_cardinality_remove_element}.
        Take $x$ such that $A\setminus\{y\}=\{x\}$ by
            \obv{finite_cardinality_one_iff_singleton}.
        $x\in A\setminus\{y\}$ by \auto{singleton_iff}.
        Hence $x\in A$ and $x\neq y$ by
            \obv{setminus,singleton_iff}.
        $\cons{y}{A\setminus\{y\}}=A$ by \obv{cons_remove}.
        Hence $A=\{x,y\}$ by
            \auto{cons_iff,singleton_iff,setext}.
        Follows by \auto.
    \end{subproof}
    Show that if there exist $x,y$ such that $x\neq y$ and
            $A=\{x,y\}$, then $A\approx\cardTwo$.
    \begin{subproof}
        Assume there exist $x,y$ such that $x\neq y$ and $A=\{x,y\}$.
        Take $x,y$ such that $x\neq y$ and $A=\{x,y\}$.
        $\{y\}$ has finite cardinality $\cardOne$ by
            \auto{finite_cardinality_one_iff_singleton}.
        $x\notin\{y\}$ by \obv{singleton_iff}.
        Hence $\cons{x}{\{y\}}$ has finite cardinality
            $\suc{\cardOne}$ by \obv{finite_cardinality_cons}.
        $\suc{\cardOne}=\cardTwo$ by \obv{cardinal_two_successor}.
        $\cons{x}{\{y\}}=A$ by
            \obv{setext,cons_iff,singleton_iff,upair_iff}.
        Hence $A$ has finite cardinality $\cardTwo$.
        Take $f$ such that $f\in\bijections{\cardTwo}{A}$ by
            \obv{finite_cardinality}.
        Hence $\cardTwo\approx A$ by \obv{equinum}.
        Follows by \obv{equinum_sym}.
    \end{subproof}
    Follows by \auto.
\end{proof}

\begin{theorem}\label{finite_iff_injects_into_natural}
    $X$ is finite iff there exists a natural number $n$ such that
        there exists an injective function from $X$ to $n$.
\end{theorem}
\begin{proof}
    Show that if $X$ is finite, then there exists a natural number $n$
        such that there exists an injective function from $X$ to $n$.
    \begin{subproof}
        Assume $X$ is finite.
        Take a natural number $n$ such that there exists
            $f\in\bijections{n}{X}$ by \auto{finite}.
        Take $f$ such that $f\in\bijections{n}{X}$ by \auto.
        $\converse{f}\in\bijections{X}{n}$
            by \obv{bijection_converse_is_bijection}.
        Hence $\converse{f}\in\funs{X}{n}$
            by \obv{bijections_to_funs}.
        $\converse{f}$ is injective by \obv{bijections}.
        Thus $\converse{f}$ is an injective function from $X$ to $n$
            by \obv{funs_is_function,injective_function}.
    \end{subproof}
    Show that if there exists a natural number $n$ such that
        there exists an injective function from $X$ to $n$, then
        $X$ is finite.
    \begin{subproof}
        Assume there exists a natural number $n$ such that
            there exists an injective function from $X$ to $n$.
        Take a natural number $n$ such that there exists an injective
            function $f$ from $X$ to $n$.
        Take $f$ such that $f$ is an injective function from $X$ to $n$.
        $f\in\funs{X}{n}$ and $f$ is injective
            by \obv{injective_function,funs_is_function}.
        Hence $f\in\inj{X}{n}$
            by \auto{inj,funs,funs_is_function,injective_function}.
        $\ran{f}\subseteq n$ by \obv{inj,funs_ran}.
        $n$ is finite by \obv{natural_number_is_finite}.
        Hence $\ran{f}$ is finite by \obv{subseteq_finite_is_finite}.
        $f\in\surj{X}{\ran{f}}$ by \obv{inj,fun_to_surj}.
        $f$ is injective
            by \obv{inj,funs,funs_is_function,injective_function}.
        Hence $f\in\bijections{X}{\ran{f}}$ by \obv{bijections}.
        Thus $X\approx\ran{f}$ by \obv{equinum}.
        $\ran{f}\approx X$ by \obv{equinum_sym}.
        Follows by \obv{finite_equinumerous}.
    \end{subproof}
    Follows by \auto.
\end{proof}

\begin{proposition}\label{cardinal_sum_natural_successor}
    Let $n$ be a natural number.
    Then $\cardSum{n}{\cardOne}=\suc{n}$.
\end{proposition}
\begin{proof}
    $n\times\{\sumLeftTag\}\approx n$
        by \obv{times_singleton_right_equinum}.
    $\cardOne\times\{\sumRightTag\}\approx\cardOne$
        by \obv{times_singleton_right_equinum}.
    $\cardOne=\{\emptyset\}$
        by \obv{cardinal_one,setext,suc_iff,singleton_iff,emptyset}.
    Hence $\cardOne\approx\{n\}$ by \auto{singletons_equinum}.
    Hence $\cardOne\times\{\sumRightTag\}\approx\{n\}$
        by \obv{equinum_tran}.
    $(n\times\{\sumLeftTag\})\inter
        (\cardOne\times\{\sumRightTag\})=\emptyset$
        by \obv{tagged_sum_components_disjoint}.
    $n\inter\{n\}=\emptyset$
        by \auto{setext,inter,singleton_iff,emptyset,in_irrefl}.
    Hence $\taggedSum{n}{\cardOne}\approx n\union\{n\}$
        by \auto{tagged_sum,disjoint_union_equinum_congruent}.
    $n\union\{n\}=\suc{n}$
        by \auto{setext,union_iff,suc_iff,singleton_iff}.
    Hence $\card{\taggedSum{n}{\cardOne}}=\card{\suc{n}}$
        by \obv{card_eq_iff_equinumerous}.
    $\suc{n}$ is a natural number by \obv{omega_successor}.
    Hence $\card{\suc{n}}=\suc{n}$
        by \obv{card_of_natural_number}.
    Follows by \obv{cardinal_sum}.
\end{proof}

\begin{theorem}\label{cardinal_sum_natural_numbers}
    Let $m,n$ be natural numbers.
    Then $\cardSum{m}{n}=m+n$.
\end{theorem}
\begin{proof}
    Let $C=\{k\in\naturals\mid
        \text{$\cardSum{m}{k}=m+k$}\}$.
    $m$ is a cardinal by \obv{natural_number_is_cardinal}.
    $\emptyset\in C$
        by
        \auto{omega_empty,cardinal_zero,cardinal_sum_zero_right,num_addition_zero_right}.
    Show that for every $k\in C$ we have $\suc{k}\in C$.
    \begin{subproof}
        Fix $k\in C$.
        Then $k$ is a natural number and
            $\cardSum{m}{k}=m+k$.
        $m+k$ is a natural number by \obv{num_addition_on_naturals}.
        $\cardSum{k}{\cardOne}=\suc{k}$
            by \obv{cardinal_sum_natural_successor}.
        $\cardSum{m}{\suc{k}}
            =\cardSum{\cardSum{m}{k}}{\cardOne}$
            by \auto{cardinal_sum_associative}.
        Hence $\cardSum{m}{\suc{k}}
            =\cardSum{m+k}{\cardOne}$.
        Hence $\cardSum{m}{\suc{k}}=\suc{m+k}$
            by \obv{cardinal_sum_natural_successor}.
        $m+\suc{k}=\suc{m+k}$
            by \obv{num_addition_successor_right}.
        $\suc{k}$ is a natural number by \obv{omega_successor}.
    \end{subproof}
    Hence $C$ is an inductive set by \obv{inductive_set}.
    Hence every natural number is an element of $C$
        by \obv{omega_minimal}.
    Thus $n\in C$.
\end{proof}

\begin{theorem}\label{cardinal_product_natural_numbers}
    Let $m,n$ be natural numbers.
    Then $\cardProduct{m}{n}=\naturalMul{m}{n}$.
\end{theorem}
\begin{proof}
    Let $C=\{k\in\naturals\mid
        \text{$\cardProduct{m}{k}=\naturalMul{m}{k}$}\}$.
    $m$ is a cardinal by \obv{natural_number_is_cardinal}.
    $\emptyset\in C$
        by
        \auto{omega_empty,cardinal_zero,cardinal_product_zero_right,natural_multiplication_zero_right}.
    Show that for every $k\in C$ we have $\suc{k}\in C$.
    \begin{subproof}
        Fix $k\in C$.
        Then $k$ is a natural number and
            $\cardProduct{m}{k}=\naturalMul{m}{k}$.
        $\naturalMul{m}{k}$ is a natural number
            by \obv{natural_multiplication_closed}.
        $\cardSum{k}{\cardOne}=\suc{k}$
            by \obv{cardinal_sum_natural_successor}.
        $\cardProduct{m}{\suc{k}}
            =\cardSum{\cardProduct{m}{k}}
                {\cardProduct{m}{\cardOne}}$
            by \auto{cardinal_product_distributes_over_sum}.
        $\cardProduct{m}{\cardOne}=m$
            by \obv{cardinal_product_one_right}.
        Hence $\cardProduct{m}{\suc{k}}
            =\cardSum{\naturalMul{m}{k}}{m}$.
        Hence $\cardProduct{m}{\suc{k}}
            =\naturalMul{m}{k}+m$
            by \obv{cardinal_sum_natural_numbers}.
        $\naturalMul{m}{\suc{k}}=\naturalMul{m}{k}+m$
            by \obv{natural_multiplication_successor_right}.
        $\suc{k}$ is a natural number by \obv{omega_successor}.
    \end{subproof}
    Hence $C$ is an inductive set by \obv{inductive_set}.
    Hence every natural number is an element of $C$
        by \obv{omega_minimal}.
    Thus $n\in C$.
\end{proof}

\begin{theorem}\label{naturals_is_infinite}
    $\naturals$ is infinite.
\end{theorem}
\begin{proof}
    Suppose not.
    $\naturals$ is finite.
    $\naturals$ is an ordinal by \obv{omega_is_ordinal}.
    Hence $\naturals$ is a natural number
        by \obv{finite_ordinal_is_natural_number}.
    Hence $\naturals\in\naturals$.
    Contradiction by \obv{in_irrefl}.
\end{proof}

\begin{theorem}\label{naturals_is_cardinal}
    $\naturals$ is a cardinal.
\end{theorem}
\begin{proof}
    $\naturals$ is an ordinal by \obv{omega_is_ordinal}.
    Show that for every $n\in\naturals$ we have
        $\lnot(n\approx\naturals)$.
    \begin{subproof}
        Fix $n\in\naturals$.
        $n$ is finite by \obv{natural_number_is_finite}.
        Suppose $n\approx\naturals$.
        Hence $\naturals$ is finite by \obv{finite_equinumerous}.
        Contradiction by \obv{naturals_is_infinite}.
    \end{subproof}
    Hence $\naturals$ is an initial ordinal by \obv{initial_ordinal}.
    Follows by \obv{cardinal}.
\end{proof}

\begin{corollary}\label{card_of_naturals}
    $\card{\naturals}=\naturals$.
\end{corollary}
\begin{proof}
    Follows by \obv{naturals_is_cardinal,card_of_cardinal}.
\end{proof}

\begin{theorem}\label{below_omega_iff_finite}
    $\cardLt{F}{\naturals}$ iff $F$ is finite.
\end{theorem}
\begin{proof}
    $\card{\naturals}=\naturals$ by \obv{card_of_naturals}.
    Show that if $\cardLt{F}{\naturals}$, then $F$ is finite.
    \begin{subproof}
        Assume $\cardLt{F}{\naturals}$.
        Hence $\card{F}\in\card{\naturals}$ by
            \obv{card_elem_iff_cardinal_lt}.
        Hence $\card{F}\in\naturals$.
        Hence $\card{F}$ is finite by
            \obv{natural_number_is_finite}.
        $\card{F}\approx F$ by \obv{card_equinumerous}.
        Follows by \obv{finite_equinumerous}.
    \end{subproof}
    Show that if $F$ is finite, then $\cardLt{F}{\naturals}$.
    \begin{subproof}
        Assume $F$ is finite.
        Take a natural number $n$ such that there exists
            $f\in\bijections{n}{F}$ by \auto{finite}.
        Take $f$ such that $f\in\bijections{n}{F}$ by \auto.
        Hence $n\approx F$ by \obv{equinum}.
        Hence $\card{n}=\card{F}$ by
            \obv{card_eq_iff_equinumerous}.
        $\card{n}=n$ by \obv{card_of_natural_number}.
        Hence $\card{F}\in\card{\naturals}$.
        Follows by \obv{card_elem_iff_cardinal_lt}.
    \end{subproof}
    Follows by \auto.
\end{proof}

\begin{theorem}\label{at_most_countable_iff_card_subseteq_naturals}
    $X$ is at most countable iff $\card{X}\subseteq\naturals$.
\end{theorem}
\begin{proof}
    $\card{\naturals}=\naturals$ by \obv{card_of_naturals}.
    Follows by
        \auto{at_most_countable,cardinal_le,card_subseteq_iff_cardinal_le}.
\end{proof}

\begin{theorem}\label{countably_infinite_iff_card_naturals}
    $X$ is countably infinite iff $\card{X}=\naturals$.
\end{theorem}
\begin{proof}
    Show that if $X$ is countably infinite, then
        $\card{X}=\naturals$.
    \begin{subproof}
        Assume $X$ is countably infinite.
        Hence $X$ is at most countable and $X$ is infinite
            by \obv{countably_infinite}.
        Hence $\card{X}\subseteq\naturals$
            by \obv{at_most_countable_iff_card_subseteq_naturals}.
        $\card{X}$ and $\naturals$ are ordinals
            by \obv{card_is_ordinal,omega_is_ordinal}.
        Hence $\card{X}\in\naturals$ or
            $\card{X}=\naturals$
            by \obv{ordinal_subseteq_implies_elem_or_eq}.
        Show that $\card{X}\notin\naturals$.
        \begin{subproof}
            Suppose $\card{X}\in\naturals$.
            Hence $\card{X}$ is finite
                by \obv{natural_number_is_finite}.
            $\card{X}\approx X$ by \obv{card_equinumerous}.
            Hence $X$ is finite by \obv{finite_equinumerous}.
            Contradiction.
        \end{subproof}
    \end{subproof}
    Show that if $\card{X}=\naturals$, then
        $X$ is countably infinite.
    \begin{subproof}
        Assume $\card{X}=\naturals$.
        $\card{X}\subseteq\naturals$ by \obv{subseteq_refl}.
        Hence $X$ is at most countable
            by \obv{at_most_countable_iff_card_subseteq_naturals}.
        Show that $X$ is infinite.
        \begin{subproof}
            Suppose not.
            Hence $X$ is finite.
            $X\approx\card{X}$ by \obv{equinumerous_card}.
            Hence $\card{X}$ is finite by \obv{finite_equinumerous}.
            Hence $\naturals$ is finite.
            Contradiction by \obv{naturals_is_infinite}.
        \end{subproof}
        Follows by \obv{countably_infinite}.
    \end{subproof}
    Follows by \auto.
\end{proof}

\begin{corollary}\label{enumerable_iff_card_subseteq_naturals}
    $X$ is enumerable iff $\card{X}\subseteq\naturals$.
\end{corollary}
\begin{proof}
    Follows by
        \obv{countable_iff_enumerable,at_most_countable_iff_card_subseteq_naturals}.
\end{proof}
